\chapter{Logische Gatter}\label{ch:logische_gatter}

\section{Bits -- die kleinste Informationseinheit}

Ein Computer verarbeitet Informationen ausschliesslich in Form von \textbf{Bits}. Ein Bit (von \emph{binary digit}, dt. Binärstelle) kann genau zwei Zustände annehmen: \texttt{0} und \texttt{1}. Physikalisch entspricht dies auf einem Chip zwei verschiedenen Spannungsniveaus -- typischerweise \enquote{keine Spannung} ($\approx 0\,\mathrm{V}$) für den Wert \texttt{0} und \enquote{hohe Spannung} ($\approx 3{,}3\,\mathrm{V}$ oder $5\,\mathrm{V}$) für den Wert \texttt{1}.

\begin{mydefinition}[Bit und Byte]
    Ein \textbf{Bit} ist die kleinste digitale Informationseinheit und nimmt einen der zwei Werte $\{0, 1\}$ an.
    
    Ein \textbf{Byte} besteht aus $8$ Bits. Mit einem Byte können $2^8 = 256$ verschiedene Werte dargestellt werden.
\end{mydefinition}

Grössere Datenmengen werden in üblichen Vielfachen angegeben:

\begin{table}[H]
    \centering
    \begin{tabular}{l|l|l}
        \textbf{Einheit} & \textbf{Abkürzung} & \textbf{Grösse} \\\midrule\midrule
        Kilobyte & KB & $10^3$ Bytes $\approx 1\,000$ Bytes \\\midrule
        Megabyte & MB & $10^6$ Bytes $\approx 1\,000\,000$ Bytes \\\midrule
        Gigabyte & GB & $10^9$ Bytes \\\midrule
        Terabyte & TB & $10^{12}$ Bytes \\
    \end{tabular}
    \caption{Gängige Einheiten für Datenmenge}
    \label{tab:einheiten}
\end{table}

\begin{myexample}[Grösse von Dateien]
    Ein kurzes Textdokument benötigt wenige Kilobytes, ein Foto einige Megabytes und ein HD-Film mehrere Gigabytes Speicherplatz.
\end{myexample}

\section{Logische Grundverknüpfungen}

Die Grundlage aller Rechenoperationen eines Computers bilden sogenannte \textbf{Boolesche Funktionen}: Funktionen, die einen oder mehrere Bits als Eingabe erhalten und einen Bit als Ausgabe liefern. Die drei elementaren Verknüpfungen sind \textbf{NOT}, \textbf{AND} und \textbf{OR}.

\subsection{NOT -- Negation}

\begin{mydefinition}[NOT-Gatter]
    Das \textbf{NOT-Gatter} (auch: Inverter) hat \emph{einen} Eingang $x$ und liefert den negierten Wert $\overline{x}$ als Ausgang. Es gilt:
    \[
        \overline{0} = 1 \qquad \overline{1} = 0.
    \]
\end{mydefinition}

\begin{figure}[H]
    \centering
    \begin{tikzpicture}[circuit logic US, scale=1.4]
        \node[not gate, draw] (not1) at (0,0) {};
        \draw (not1.input) -- ++(-0.7,0) node[left] {$x$};
        \draw (not1.output) -- ++(0.7,0) node[right] {$\overline{x}$};
    \end{tikzpicture}
    \caption{Symbol des NOT-Gatters.}
    \label{fig:not_gatter}
\end{figure}

\begin{table}[H]
    \centering
    \begin{tabular}{c||c}
        $x$ & $\overline{x}$ \\\midrule\midrule
        0 & 1 \\\midrule
        1 & 0 \\
    \end{tabular}
    \caption{Wahrheitstabelle des NOT-Gatters.}
    \label{tab:not}
\end{table}

\subsection{AND -- Konjunktion}

\begin{mydefinition}[AND-Gatter]
    Das \textbf{AND-Gatter} hat \emph{zwei} Eingänge $x$ und $y$ und liefert genau dann den Ausgang \texttt{1}, wenn \emph{beide} Eingänge \texttt{1} sind:
    \[
        x \wedge y = 1 \iff x = 1 \text{ und } y = 1.
    \]
\end{mydefinition}

\begin{figure}[H]
    \centering
    \begin{tikzpicture}[circuit logic US, scale=1.4]
        \node[and gate, draw] (and1) at (0,0) {};
        \draw (and1.input 1) -- ++(-0.7,0) node[left] {$x$};
        \draw (and1.input 2) -- ++(-0.7,0) node[left] {$y$};
        \draw (and1.output) -- ++(0.7,0) node[right] {$x \wedge y$};
    \end{tikzpicture}
    \caption{Symbol des AND-Gatters.}
    \label{fig:and_gatter}
\end{figure}

\begin{table}[H]
    \centering
    \begin{tabular}{c|c||c}
        $x$ & $y$ & $x \wedge y$ \\\midrule\midrule
        0 & 0 & 0 \\\midrule
        0 & 1 & 0 \\\midrule
        1 & 0 & 0 \\\midrule
        1 & 1 & 1 \\
    \end{tabular}
    \caption{Wahrheitstabelle des AND-Gatters.}
    \label{tab:and}
\end{table}

\begin{myexample}[Türschloss mit zwei Schlüsseln]
    Stellen Sie sich ein Schloss vor, das sich nur öffnet, wenn \emph{zwei} Schlüssel gleichzeitig umgedreht werden. Dieses Schloss verhält sich wie ein AND-Gatter: Die Tür öffnet sich (\texttt{1}) nur dann, wenn Schlüssel 1 \emph{und} Schlüssel 2 gedreht sind.
\end{myexample}

\subsection{OR -- Disjunktion}

\begin{mydefinition}[OR-Gatter]
    Das \textbf{OR-Gatter} hat \emph{zwei} Eingänge $x$ und $y$ und liefert genau dann den Ausgang \texttt{1}, wenn \emph{mindestens einer} der Eingänge \texttt{1} ist:
    \[
        x \vee y = 1 \iff x = 1 \text{ oder } y = 1.
    \]
\end{mydefinition}

\begin{figure}[H]
    \centering
    \begin{tikzpicture}[circuit logic US, scale=1.4]
        \node[or gate, draw] (or1) at (0,0) {};
        \draw (or1.input 1) -- ++(-0.7,0) node[left] {$x$};
        \draw (or1.input 2) -- ++(-0.7,0) node[left] {$y$};
        \draw (or1.output) -- ++(0.7,0) node[right] {$x \vee y$};
    \end{tikzpicture}
    \caption{Symbol des OR-Gatters.}
    \label{fig:or_gatter}
\end{figure}

\begin{table}[H]
    \centering
    \begin{tabular}{c|c||c}
        $x$ & $y$ & $x \vee y$ \\\midrule\midrule
        0 & 0 & 0 \\\midrule
        0 & 1 & 1 \\\midrule
        1 & 0 & 1 \\\midrule
        1 & 1 & 1 \\
    \end{tabular}
    \caption{Wahrheitstabelle des OR-Gatters.}
    \label{tab:or}
\end{table}

\section{Das NAND-Gatter: Ein universelles Gatter}\label{sec:nand}

\begin{mydefinition}[NAND-Gatter]
    Das \textbf{NAND-Gatter} (von \emph{Not AND}) liefert die Negation des AND:
    \[
        x \uparrow y \;=\; \overline{x \wedge y}.
    \]
    Es gibt \texttt{0} aus, wenn \emph{beide} Eingänge \texttt{1} sind; sonst immer \texttt{1}.
\end{mydefinition}

\begin{figure}[H]
    \centering
    \begin{tikzpicture}[circuit logic US, scale=1.4]
        \node[nand gate, draw] (nand1) at (0,0) {};
        \draw (nand1.input 1) -- ++(-0.7,0) node[left] {$x$};
        \draw (nand1.input 2) -- ++(-0.7,0) node[left] {$y$};
        \draw (nand1.output) -- ++(0.7,0) node[right] {$\overline{x \wedge y}$};
    \end{tikzpicture}
    \caption{Symbol des NAND-Gatters (AND mit Kreis am Ausgang).}
    \label{fig:nand_gatter}
\end{figure}

\begin{table}[H]
    \centering
    \begin{tabular}{c|c||c}
        $x$ & $y$ & $\overline{x \wedge y}$ \\\midrule\midrule
        0 & 0 & 1 \\\midrule
        0 & 1 & 1 \\\midrule
        1 & 0 & 1 \\\midrule
        1 & 1 & 0 \\
    \end{tabular}
    \caption{Wahrheitstabelle des NAND-Gatters.}
    \label{tab:nand}
\end{table}

Das NAND-Gatter ist besonders bedeutsam: Durch Kombinationen von NAND-Gattern lassen sich \emph{alle} anderen logischen Funktionen (NOT, AND, OR, \ldots) nachbauen. Man sagt, NAND ist \textbf{funktional vollständig}. Tatsächlich sind moderne Computerchips intern zu einem grossen Teil aus NAND-Gattern aufgebaut -- Milliarden davon auf der Fläche eines Fingernagels.

\begin{myexercise}
    Zeigen Sie, dass sich das NOT-Gatter allein mit NAND-Gattern realisieren lässt. Stellen Sie dazu die Wahrheitstabelle von $x \uparrow x$ (also: $x$ NAND $x$) auf und vergleichen Sie diese mit der Wahrheitstabelle des NOT-Gatters.
    \begin{myanswer}
        \begin{center}
            \begin{tabular}{c||c|c}
                $x$ & $x \wedge x$ & $\overline{x \wedge x}$ \\\midrule\midrule
                0 & 0 & 1 \\\midrule
                1 & 1 & 0 \\
            \end{tabular}
        \end{center}
        Die Spalte $\overline{x \wedge x}$ ist identisch mit der Wahrheitstabelle von NOT($x$). Ein NOT-Gatter kann also durch ein einziges NAND-Gatter, bei dem beide Eingänge verbunden sind, realisiert werden.
    \end{myanswer}
\end{myexercise}

\section{Kombinierte Schaltungen: der Halbaddierer}\label{sec:halbaddierer}

Logische Gatter werden zu komplexeren Schaltungen kombiniert. Als konkretes Beispiel betrachten wir die binäre Addition zweier einstelliger Zahlen. Wenn zwei Bits $x$ und $y$ addiert werden, entsteht eine zweistellige Summe:

\begin{center}
    \begin{tabular}{cc}
        $x$ & $y$ \\\midrule
        0 & 0 \\
        0 & 1 \\
        1 & 0 \\
        1 & 1 \\
    \end{tabular}
    \quad$\longrightarrow$\quad
    \begin{tabular}{cc}
        Übertrag $c$ & Summe $s$ \\\midrule
        0 & 0 \\
        0 & 1 \\
        0 & 1 \\
        1 & 0 \\
    \end{tabular}
\end{center}

Die Summe $s$ entspricht $x \oplus y$ (XOR: genau dann \texttt{1}, wenn die Bits verschieden sind), der Übertrag $c$ entspricht $x \wedge y$ (AND):

\begin{figure}[H]
    \centering
    \begin{tikzpicture}[circuit logic US, scale=1.3]
        % XOR gate for sum
        \node[xor gate, draw] (xor1) at (2, 0.5) {};
        % AND gate for carry
        \node[and gate, draw] (and1) at (2,-0.9) {};

        % Inputs x
        \draw (xor1.input 1) -- ++(-0.5,0) coordinate (xinxor) -- ++(-0.7,0) node[left] {$x$};
        \draw (and1.input 1) -- ++(-0.5,0) coordinate (xinand);
        \draw (xinxor) -- ++(0,-0.5) -- (xinand);

        % Inputs y
        \draw (xor1.input 2) -- ++(-0.3,0) coordinate (yinxor) -- ++(-0.9,0) node[left] {$y$};
        \draw (and1.input 2) -- ++(-0.3,0) coordinate (yinand);
        \draw (yinxor) -- ++(0,-0.5) -- (yinand);

        % Outputs
        \draw (xor1.output) -- ++(0.8,0) node[right] {$s$ (Summe)};
        \draw (and1.output) -- ++(0.8,0) node[right] {$c$ (Übertrag)};
    \end{tikzpicture}
    \caption{Schaltplan eines Halbaddierers: XOR liefert die Summe, AND den Übertrag.}
    \label{fig:halbaddierer}
\end{figure}

Solche Addierwerke werden in einer echten CPU zu vollständigen \textbf{Addierwerken} (englisch: \emph{adder}) zusammengeschaltet, die beliebig grosse Ganzzahlen addieren können.

\begin{myexercise}
    Überprüfen Sie anhand der Wahrheitstabelle, dass der Halbaddierer die binäre Addition korrekt berechnet: $1 + 1 = 10_2$ (dezimal: $2$), also Summe $= 0$ und Übertrag $= 1$.
    \begin{myanswer}
        Für $x=1, y=1$: XOR$(1,1) = 0$ (Summe), AND$(1,1) = 1$ (Übertrag). Das Ergebnis ist also die zweistellige Binärzahl \texttt{10}, was $2$ im Dezimalsystem entspricht. Korrekt.
    \end{myanswer}
\end{myexercise}
